Pairs of firms convicted of the same procurement cartel exchange, on average, 74 workers between them over eight years. Pairs of non-cartel firms drawn from the same bidder universe exchange essentially none. This paper asks whether pair-level worker flow, linked to product-market participation, can be built into a usable cartel screen. We link six CADE-convicted S\~{a}o Paulo cartels---33 firms, 100 within-cartel pairs---to the universe of Brazilian formal employment (RAIS) and the bid-level electronic procurement platform BEC. A class-balanced logistic composite of log pair-level worker flow and log count of items on which both firms bid achieves a leave-one-CADE-cartel-out macro AUC of $0.845$ and a pooled AUC of $0.809$. The composite dominates three trivial baselines (four-digit industry match, common-auction indicator, common-buyer indicator), the overlap-only reference, and the static labor-cosine specification used in an earlier version of this project. A label-permutation test on the pooled AUC rejects the no-signal null at the 2\% level. A cluster bootstrap over the six cartels places the macro AUC in a 95\% confidence interval of $[0.66, 0.96]$; the bootstrap 95\% CI on the marginal contribution of worker flow over product-market overlap alone is $[-0.07, +0.27]$ with empirical $p \approx 0.10$, reflecting the precision limit of a five-effective-cartel validation design. The classical variance-based screen of Abrantes-Metz et al.\ (2006) fails entirely on our preg\~{a}o-auction data, which we attribute to the descending-auction format mechanically compressing within-auction bid distributions; the proposed fingerprints are the practical alternative for dynamic procurement auctions. We characterize failure modes transparently: small, short-lived cartels leave detectable traces in neither labor nor product-market dimensions, giving authorities a defensible operating envelope. The framework uses only a matched employer--employee register and bidder identities, and can be calibrated on local ground truth in any jurisdiction with these two data layers.
